\documentclass[12pt]{article}
\usepackage[utf8]{inputenc}
\usepackage[T1]{fontenc}
\usepackage{hyperref}
\title{Ulau-Suain Language Resources for K-12 Curriculum Development}
\author{Compiled by Claude Sonnet 5.5 (claude-sonnet-5-5) for Andr\'as Kornai}
\date{}
\begin{document}
\maketitle

\section{Introduction}
Ulau-Suain (ISO 639-3: svb; Glottocode ulau1237) is an Austronesian (Oceanic, Schouten) language spoken on the Sepik coast of Sandaun Province, Papua New Guinea, in the Ulau and Suain wards of East Aitape Rural LLG. Sources give about 2,800 speakers (an older figure); Joshua Project lists a higher number. Very little material exists. The search found only reference, archival and mission-audio items, and no school materials.

\section{Primary Educational Resources}
No primers, textbooks or school programmes for Ulau-Suain were found in the search results. The language is not known to be taught in schools.

\section{Linguistic and Reference Materials}
\begin{itemize}
\item Wikipedia, Ulau-Suain language: \url{https://en.wikipedia.org/wiki/Ulau-Suain_language}
\item Ulau Suain numerals (Max Planck Institute for Evolutionary Anthropology, Numeral Systems of Chan and Wuchen): \url{https://lingweb.eva.mpg.de/channumerals/Ulau-Suain.htm}
\item Endangered Languages Project entry: \url{https://endangeredlanguages.com/lang/10523?hl=fr}
\end{itemize}
The search results mention that Paradisec holds Ulau-Suain materials from Arthur Capell and Malcolm Ross, but no direct URL appeared in the results.

\section{Religious and Community Resources}
\begin{itemize}
\item Joshua Project, Suain people group profile (reports that no Bible translation is available): \url{https://joshuaproject.net/index.php/people_groups/15096}
\item Joshua Project language page for svb: \url{https://www.joshuaproject.net/index.php/languages/svb}
\item Scripture Earth listing for svb: \url{https://www.scriptureearth.org/00eng.php?iso=svb}
\end{itemize}

\section{Digital and Audio Resources}
\begin{itemize}
\item Global Recordings Network, Suain: Ulau language (audio Bible stories and evangelistic recordings): \url{https://globalrecordings.net/en/language/587}
\end{itemize}

\section{Curriculum Development Context}
There is almost no published material for Ulau-Suain, so a curriculum would have to be built from community knowledge, with archival and linguistic documentation as support. Nothing found indicates a standard orthography or school use. The speaker estimates conflict, and the Ethnologue data cited by the sources is old. Any programme would need community-led documentation first.

\end{document}